% Generated draft: requires mathematical and source review.
\documentclass[11pt]{article}
\usepackage[margin=25mm]{geometry}
\usepackage{fontspec}
\setmainfont{Latin Modern Roman}
\usepackage{amsmath,amssymb,mathtools,unicode-math}
\setmathfont{Latin Modern Math}
\usepackage{enumitem,needspace,xurl,hyperref}
\setlength{\parindent}{0pt}
\setlength{\parskip}{5pt}
\setlength{\emergencystretch}{3em}
\newcommand{\N}{\mathbb{N}}
\newcommand{\Z}{\mathbb{Z}}
\newcommand{\Q}{\mathbb{Q}}
\newcommand{\R}{\mathbb{R}}
\newcommand{\C}{\mathbb{C}}
\newcommand{\F}{\mathbb{F}}
\newcommand{\A}{\mathbb{A}}
\newcommand{\PP}{\mathbb{P}}
\newcommand{\E}{\mathbb{E}}
\newcommand{\eps}{\varepsilon}
\newcommand{\dd}{\,\mathrm{d}}
\begin{document}
\section*{UnsolvedMath ID: 30001110}
\textbf{Source problem code:} catalogue entry
\subsection{Definitions and mathematical statement}
Let \(p=3\), and let \(\pi:C\to\PP^1\) be the wildly ramified Artin--Schreier cover intended in the supplied source, with exactly one branch point, namely \(\infty\). An Artin--Schreier cover in characteristic \(3\) is a cyclic cover of degree \(3\) described at the function-field level by an equation \(y^3-y=f\). The ground field and the particular cover are not specified in the record.

Write \(R_\pi\) for the ring called the universal deformation ring of \(\pi\) in the source. The supplied record asserts a presentation
\[
R_\pi=R_\infty\cong W[[t_1,t_2]]/I,\qquad I\ne(0),
\]
where \(W[[t_1,t_2]]\) denotes the formal power series ring in two variables over the coefficient ring denoted by \(W\). The record does not define \(W\), the deformation functor, or the local deformation problem represented by \(R_\infty\); the displayed presentation is therefore retained as a source-supplied premise for the intended example, rather than asserted for every Artin--Schreier cover with one branch point.

\textbf{Question.} For this intended deformation problem and presentation, is \(I\) principal? Explicitly, setting \(A=W[[t_1,t_2]]\), does there exist \(F\in A\setminus\{0\}\) such that
\[
I=(F)=\{FG:G\in A\}?
\]
The missing specification of the cover, coefficient ring, and deformation problem must be recovered before this can be treated as a fully determined mathematical question.
\subsection{Short English statement}
For the intended Artin--Schreier cover in characteristic \(3\), branched only at infinity, the supplied source describes its deformation ring as \(W[[t_1,t_2]]/I\) with \(I\ne(0)\). The question asks whether this ideal can be generated by one element. The record leaves the particular cover, the coefficient ring \(W\), and the deformation conditions unspecified.
\subsection{Sources}
\begin{itemize}
\item Supplied local UnsolvedMath registry record, permanent ID \(30001110\), title \emph{Principality of Artin--Schreier Deformation Ideals}, source problem code \texttt{catalogue entry}, in the Oberwolfach Reports Open Problems collection. No individual proposer is supplied; the proposed year is recorded as 2008.
\item Supplied report citation: \emph{Mini-Workshop: Group Actions on Curves: Reduction and Lifting} (2009), Oberwolfach Reports, pp.~3029--3064, DOI \texttt{10.4171/owr/2008/54}. Supplied source URL and literature-source link labeled ``original OWR report'': \url{https://doi.org/10.4171/owr/2008/54}. The report text was not independently inspected for this draft.
\item Literature assessment dated \texttt{2026-08-21}: this is supplied dated curation, not an independently verified assessment. It labels the problem open and states that the report leaves this principality question open. It further reports later structural results for local deformation functors and that no resolution of this specific question was found. No separate references for those later results are supplied. The background explicitly requests independent verification of time-sensitive status and cited claims.
\item The registry labels the statement \texttt{corrected\_verified}; its accompanying verification note describes preservation of the statement with normalization of malformed or undelimited notation. This supplied label does not constitute mathematical verification of the present draft.
\end{itemize}
\end{document}
